% ============================================================
% HEAS — WSC 2026 Submission  (Version 7 / v0.2)
% Track: Modeling Methodology
%
% INSTRUCTIONS FOR COMPILATION:
%   Run: pdflatex paper && bibtex paper && pdflatex paper && pdflatex paper
%   (wscpaperproc.cls and support files are in this directory)
%
% EDITORIAL CHANGES APPLIED (v0.2 / Version 7):
%   [TASK 1] DONE: Removed Figure 2 (fig:champion32) from §5.4
%   [TASK 2] DONE: Rewrote §5.4 closing sentence with OOD decoupling premise
%   [TASK 3] DONE: Split aggregation inconsistency into structural + outcome levels
%   [TASK 4] DONE: Clarified binomial test directionality (H_0, H_a)
%   [TASK 5] DONE: Reframed d=0.17 interpretation (tight distributions, decoupling)
%   [TASK 7+8] DONE: Tightened §3.5 Design Constraints (added aggregation clarification, condensed NSGA-II sentence)
%   [TASK 9] DONE: Reframed §4 Case Studies opening (composition vehicles, decoupling validation)
%   [TASK 10] DONE: Tightened §1.2 Contributions (impossibility for aggregation to differ, semantic correctness)
%   [TASK 11] DONE: Verified Appendix table (S6b confirmed for 32-scenario, no fig:champion32 refs)
%   [TASK 12] DONE: Polished Table 2 caption (added confirmation of metric contract necessity)
%   [NEW TASK A] DONE: Added causal loop closure in §5.1 (design hypothesis validation)
%   [NEW TASK B] DONE: Strengthened Related Work positioning (architectural vs syntactic distinction)
%
% WHAT IS NEW IN V3 vs V2
%   - Abstract: "Three case studies", 1,000-step scale, 32-scenario OOD
%   - §4 Case studies: forward-reference hooks to large-scale results
%   - §5 Evaluation: restructured 4 RQs → 5 RQs
%       RQ2 upgraded: adds large-scale 20-run study (d=1.39, p<10^-6)
%       RQ3 upgraded: adds 32-scenario OOD champion robustness (32/32 wins)
%       RQ4 upgraded: adds scale heatmap + Wolf-Sheep 30-run study
%       RQ5 new:      cross-domain portability (Eco + Enterprise at scale)
%   - 4 new figures: large-scale showdown, noise stability,
%                    champion 32-scenario, scale heatmap
%   - Appendix: 7 new rows for large-scale and OOD results
%   - Conclusion: upgrade summary to match richer evidence base
%
% REQUIREMENT CHECKLIST (WSC 2026):
%   [x] Abstract <= 150 words, single paragraph, no keywords, no math symbols
%   [x] Title bold ALL CAPS (handled by wscpaperproc class)
%   [x] Section titles numbered, bold, ALL CAPS
%   [x] Times New Roman 11pt (handled by wscpaperproc class)
%   [x] 5--12 pages
%   [x] Author biographies mandatory
%   [x] No footnotes
%   [x] No page numbers (handled by wscpaperproc class)
%   [x] Verify final page count after compile  (12 pages, 2026-03-26)
%   [x] Author names/affiliations/emails confirmed camera-ready
% ============================================================

\documentclass{wscpaperproc}
\usepackage{latexsym}      % extra symbols (WSC template standard)
\usepackage{mathptmx}      % Times New Roman for math (WSC template standard)
\usepackage[T1]{fontenc}   % proper font encoding (WSC template standard)
\usepackage{booktabs}      % professional table rules
\usepackage{listings}      % code listings
\usepackage{microtype}     % better text justification
\usepackage{url}           % URL formatting
\usepackage{amsmath}       % math
\usepackage{graphicx}      % figures
\usepackage{pifont}        % \cmark / \xmark
\usepackage{multirow}      % table cell merging
\usepackage[pdftex,colorlinks=true,urlcolor=blue,citecolor=black,anchorcolor=black,linkcolor=black]{hyperref}

\graphicspath{{figs/}}

% ---- Code listing style ---
\lstset{
  basicstyle=\ttfamily\small,
  columns=flexible,
  keepspaces=true,
  frame=single,
  framerule=0.4pt,
  xleftmargin=6pt,
  xrightmargin=6pt,
  breaklines=true,
  showstringspaces=false,
  numbers=none
}

\newcommand{\cmark}{\ding{51}}
\newcommand{\xmark}{\ding{55}}

% ============================================================
\begin{document}

\WSCpagesetup{ZHANG, NIE, and ZHAO}

\title{HEAS: A HIERARCHICAL EVOLUTIONARY AGENT-BASED SIMULATION FRAMEWORK
       FOR MULTI-OBJECTIVE POLICY SEARCH}

\author{\begin{center}Ruiyu ZHANG\textsuperscript{1}, Lin NIE\textsuperscript{2}, and Xin ZHAO\textsuperscript{2}\\
[11pt]
\textsuperscript{1}Dept.~of Politics and Public Administration, The University of Hong Kong, Hong Kong SAR, CHINA\\
\textsuperscript{2}Dept.~of Applied Social Sciences, The Hong Kong Polytechnic University, Hong Kong SAR, CHINA\end{center}
}

\maketitle

\vspace{-12pt}

% ============================================================
\begin{abstract}
HEAS is a reusable framework for \emph{ODE-style and tabular} simulation-based
policy search. It addresses a recurring integration problem in this model
family: the same outcome metric is often recomputed differently across
optimization, policy comparison, and statistical validation. HEAS resolves this
through a uniform \texttt{metrics\_episode()} contract shared by all three
stages, together with a Layer/Stream/Arena composition model, a pluggable
optimizer interface, and a tournament primitive for multi-scenario ranking.
Against a Mesa~3.3.1 baseline, HEAS reduces coupling code from 160 lines to 5,
and from 42 to 5 against a reusable utility-library baseline. Three case
studies---ecological, enterprise, and a Wolf-Sheep mean-field reimplementation
of a published Mesa model---show code reuse across the targeted model family.
In the tested ecological scenarios, the evolved champion wins all 32 held-out
cases and all 24 single-anchor OOD cases.
\end{abstract}

% ============================================================
\section{INTRODUCTION}

Agent-based models are foundational for policy design across ecology,
economics, and organizational science, yet coupling an ODE-style or tabular ABM
to evolutionary optimization requires integration infrastructure that
general-purpose frameworks do not provide. This paper addresses a specific
integration failure: when researchers link an ABM to an evolutionary algorithm,
the same outcome metric is recomputed in incompatible ways across three separate
code paths---optimization, tournament evaluation, and statistical
validation---introducing silent divergence that is difficult to audit. Current
practice resolves this through bespoke coupling code (roughly 160 lines in a
representative Mesa workflow) that entangles metric definition with solver,
comparison, and inference logic.

We introduce HEAS, a framework that standardizes metric definition through a
uniform \texttt{metrics\_episode()} contract consumed identically by the
evolutionary optimizer, tournament scorer, and bootstrap confidence-interval
pipeline. This approach reduces coupling code to five lines and eliminates
conditions for metric divergence.

HEAS targets a specific niche---coupled ODE/tabular dynamics, multi-objective
evolutionary search, and multi-run validation---rather than the full ABM space
including spatially-explicit heterogeneous-agent systems. The framework's four
core components (C1--C4, Section~2) are independently useful and may be adopted
in isolation.

\subsection{The Coupling Code Problem}

Consider a researcher using Mesa to study regulatory policy in a multi-firm
economy. She wants to (1)~evolve a Pareto-optimal tax regime using NSGA-II,
(2)~compare the evolved regime against a reference policy across independent
economic scenarios using tournament voting, and (3)~report bootstrap confidence
intervals. Each step reads the same welfare metric, yet in Mesa each is an
independent code path: the DEAP fitness function extracts welfare from the
model's \texttt{DataCollector} after each run; the tournament scorer extracts
welfare from a separate episode-evaluation function, possibly with different
aggregation; and the bootstrap CI loop extracts welfare from a third data
collection loop, potentially inconsistent with both.

In practice, silent divergence is difficult to detect. Consider an experiment
where the tournament scorer reads the final step's biomass while the
evolutionary algorithm reads the episode mean biomass. The evolved champion may
become the policy with volatile-but-high peaks rather than sustainable mean
production---a consistently applied but undetected inconsistency that could
persist through publication. Such divergences are common precisely because
\emph{no single code path} connects metric definition to all three consumers.

When the researcher adds a second objective---welfare inequality---she must edit
the DataCollector reporter, the DEAP fitness function, the tournament scorer, and
the CI loop, keeping all four in sync manually. A silent divergence can persist
undetected for weeks and invalidate a publication. This pattern, which we call
\emph{coupling code}, appears in representative ABM-based policy-search
workflows~\cite{pangallo2019best,capri2023transport}.

\subsection{Contributions}
\label{sec:intro-contributions}

This paper's primary contribution is a single-source-of-truth metric framework that prevents
aggregation inconsistency in multi-objective agent-based simulation optimization.
The pattern enforces a uniform \texttt{metrics\_episode()} scalar-dict interface,
consumed identically by the evolutionary optimizer, tournament scorer, and statistical
inference pipeline. By fixing the metric contract at the framework level, HEAS
eliminates silent metric divergence by enforcing a single implementation shared across all
pipeline stages---not through explicit validation or post-hoc reconciliation, but by making
it impossible for aggregation to differ between them. Semantic correctness---whether \texttt{metrics\_episode()} correctly captures domain intent---remains the modeller's responsibility, not the framework's.

To realize this pattern, HEAS provides four implementation components:
\textbf{C1} defines the contract (\texttt{metrics\_episode()} returning
\texttt{dict[str, float]}); \textbf{C2} (Layer/Stream/Arena composition) enables
models to produce C1 without rewriting dynamics; \textbf{C3} (pluggable multi-objective
optimizer wrapper) and \textbf{C4} (tournament with four voting rules) consume C1
without internal variation. These components are independently replaceable, but none
are separate contributions---they are infrastructure for enforcing the contract.

The methodological contribution is twofold: (a)~a replicable framework design that
reduces per-project coupling code from 160 LOC to 5 LOC (97\% reduction against Mesa)
and eliminates metric divergence risk through single-source-of-truth metric enforcement, and (b)~a validated reference
implementation with empirical evidence (controlled aggregation experiment with stochastic robustness test, \S5.2;
cross-domain tournament validation, \S5.5; OOD robustness tests, \S5.5) that the
contract mechanism prevents rank reversals and supports multi-objective ABM workflows
at scale. This combination of architectural design and validated effectiveness provides
both a learning artifact and an extensible platform for simulation-based policy search.

% ============================================================
\section{RELATED WORK}

Agent-based modeling is widely used in ecology and economics, yet frameworks
like Mesa do not natively support coupled evolutionary optimization and
statistical comparison~\cite{kazil2020mesa}. Simulation optimization recognizes
this integration gap~\cite{fu2002optimization}, but search libraries (DEAP,
Optuna) and ranking methods do not standardize metric sharing across
optimization, tournament, and inference stages~\cite{fortin2012deap,akiba2019optuna}.


\begin{table}[ht]
\centering\small
\caption{Metric Contract Consistency Across ABM Policy-Search Tools.
HEAS enforces a uniform scalar-dict interface (\texttt{dict[str, float]})
consumed identically at three stages (optimization, tournament, inference).
Competitor tools address each stage independently without standardizing
the metric pathway. A checkmark (\cmark) indicates the tool provides
native support for that stage; an X (\textbf{x}) indicates users must implement
metric passing manually. Only HEAS combines all three within a unified contract.}
\label{tab:contracts}
\begin{tabular}{lccccc}
\toprule
Tool & Optimizer & Tournament & Inference & Unified Contract \\
\midrule
DEAP & \cmark & --- & --- & \textbf{x} (users write) \\
Optuna & \cmark & --- & --- & \textbf{x} (users write) \\
EMAworkbench & \cmark & \cmark & --- & \textbf{x} (per-model) \\
Mesa~3.3 & --- & --- & --- & \textbf{x} (coupling code) \\
Ax / BoTorch & \cmark & --- & --- & \textbf{x} (external) \\
HEAS & \cmark & \cmark & \cmark & \cmark \\
\bottomrule
\end{tabular}
\end{table}

Table~\ref{tab:contracts} summarizes the metric contract gap across incumbent tools.
Optimizer libraries (DEAP, Optuna) provide flexible search backends but do not
standardize how metrics are aggregated for tournament ranking or statistical inference.
EMAworkbench extends ensemble management to tournament evaluation but leaves the metric
interface unspecified---users define metrics per model. HEAS closes this gap by enforcing
a single \texttt{metrics\_episode()} contract at the framework level, eliminating
ad-hoc aggregation and making metric reuse auditable without per-project coupling code.

HEAS is positioned at this interface. It does not propose a new optimizer.
Instead, it standardizes the metric pathway shared by optimization, tournament
scoring, and inference. In that sense it is adjacent to optimizer wrappers,
ensemble managers, and reproducibility tooling, but narrower in scope and more
specific in its methodological target.

Unlike per-model utility wrappers (EMAworkbench, Mesa \texttt{DataCollector}), which require users to specify aggregation logic independently for each new domain, HEAS enforces a uniform contract at the framework level. The difference is architectural rather than syntactic: framework-level enforcement means new projects inherit consistent aggregation without writing or auditing per-project glue code.

% ============================================================
\section{THE HEAS FRAMEWORK}

HEAS is implemented in Python~3.11 and depends on DEAP for evolutionary
algorithms, NumPy for numerical computation, and SciPy for statistical
utilities. All components---EA algorithm, voting rule, statistical
estimator---can be replaced without changing simulation or metric contract code.

\subsection{Layered Composition}

This subsection defines the compositional boundary that lets models be reused
without rewriting evaluation code. The fundamental abstraction is a
\texttt{Layer} with four methods
(\texttt{step}, \texttt{metrics\_step}, \texttt{metrics\_episode},
\texttt{reset}). All downstream components read the same
\texttt{metrics\_episode()} keys.

Layers compose into \textbf{Streams}---named groups forming a coherent
sub-simulation---and Streams compose into an \textbf{Arena}, which orchestrates
step ordering, inter-stream signal routing, and episode-level aggregation. The
naming convention \texttt{stream\_name.metric\_name} ensures globally unambiguous
keys regardless of Arena size. Explicit composition boundaries make components
testable and replaceable without rewriting optimization or tournament code;
the uniform metric contract (C2) remains stable as model internals evolve.

\subsection{Native Multi-Objective EA Integration}

This subsection shows how the metric contract is consumed by optimization
without introducing a separate evaluation interface. A researcher adds
NSGA-II to a new simulation by writing only the objective function (typically
5--10 lines reading from \texttt{run\_many()} results). The EA algorithm is
replaceable through the \texttt{run\_ea()} interface. \texttt{run\_many()}
executes episodes in parallel via \texttt{ProcessPoolExecutor}, with per-episode
seeds derived deterministically for reproducibility.

\subsection{Tournament Evaluation}

This subsection defines the comparison layer that uses the same metric contract
for multi-scenario policy ranking. Comparing candidate policies is formalized as a \texttt{Tournament}: a
cross-product of participants and scenarios evaluated for $n_\text{ep}$
episodes and aggregated by argmax, majority, Borda, or Copeland rules. The
rules are intended to reveal whether rankings are robust to aggregation choice.
Universal agreement (as in Section~5.4) indicates clear dominance;
disagreement diagnoses the competitive structure.
HEAS reports rule agreement, sample complexity $P(\text{correct})$ versus
$n_\text{ep}$, and Kendall's~$\tau$~\cite{kendall1938new} stability versus
score noise.

\subsection{Statistical Infrastructure}

This subsection makes the inference layer explicit: the same episode metrics
used by optimization and tournaments also feed uncertainty quantification.
\texttt{heas/utils/stats.py} provides: \texttt{summarize\_runs()} (mean, std,
median, BCa bootstrap 95\% CI~\cite{efron1994bootstrap}),
\texttt{wilcoxon\_test()}, \texttt{cohens\_d()}, and \texttt{kendall\_tau()}
with bootstrap CI. \texttt{heas/utils/pareto.py} provides
\texttt{hypervolume()}~\cite{zitzler1998hypervolume} and
\texttt{auto\_reference\_point()}. All utilities accept \texttt{list[float]}
and return structured dicts, requiring no configuration.

\subsection{Design Constraints and Trade-offs}
\label{sec:design-constraints}

The uniform metric contract \texttt{metrics\_episode()} deliberately implements
a minimalist interface: \texttt{dict[str, float]}. This design choice reflects a
composability-versus-expressiveness trade-off.

\textbf{Minimalism enables composition.} A richer contract---e.g.,
\texttt{dict[str, list[float]]} permitting per-step time series, or
\texttt{dict[str, dict[str, float]]} enabling agent-level distributions---would
require negotiation between model internals and the optimization, tournament, and
inference layers. By constraining the contract to scalar floats, HEAS avoids
entangling aggregation logic with metric definition. A model that produces
agent-level data aggregates once at the \texttt{metrics\_episode()} boundary, then
reuses that scalar summary everywhere downstream. Richer contracts would require
per-stage aggregation choices, reintroducing the coupling-code problem HEAS targets.


\textbf{What is lost.} The scalar contract does not directly represent time-series
outputs or agent-level heterogeneity (e.g., the distribution of firm sizes, not
just mean firm profit). The contract enforces consistent aggregation \emph{of} agent data rather than removing it: each model retains full internal state, and \texttt{metrics\_episode()} may compute any aggregate statistic before returning the scalar summary. For policy-ranking purposes---the intended use case---this
loss is bounded: selecting among regulatory regimes requires aggregate summaries
as the basis for comparison, even if within-regime agent-level structure remains
unmeasured.

\textbf{NSGA-II as reference implementation.} Section~3.2 adopts NSGA-II as the
default multi-objective evolutionary algorithm based on three criteria:
(1)~established performance on policy-search benchmarks in the simulation
optimization literature; (2)~well-understood parametrization at WSC; and
(3)~ease of replacement via the pluggable \texttt{run\_ea()} interface. The
framework does not depend on this choice: \S\ref{sec:scale} shows optimizer performance is landscape-dependent, and the pluggable interface enables this comparison without code changes.

% ============================================================
\section{CASE STUDIES}

The three case studies below serve as composition vehicles to validate whether HEAS's decoupling architecture holds across distinct domain logics. The question is not whether these domains are realistically calibrated, but whether heterogeneous simulation code composes without introducing hidden coupling. Each study uses an independent \texttt{metrics\_episode()} implementation; the framework code is shared and unchanged across all three.

\subsection{Ecological Population Management}

A predator-prey Arena with a 2-gene policy (\texttt{risk},
\texttt{dispersal}) and objectives \texttt{-agg.mean\_biomass},
\texttt{agg.cv}. Five streams model climate forcing, landscape quality, prey
logistic growth, predator dynamics, and cross-stream aggregation; eight
fragmentation$\times$shock scenarios form the evaluation grid.
Section~5.3 validates this model at two simulation scales; the evolved champion
is subjected to a 32-scenario out-of-distribution robustness test in
Section~5.4.

\subsection{Enterprise Regulatory Design}

A four-layer regulatory Arena with a 4-gene policy
(\texttt{tax\_rate}, \texttt{audit\_intensity}, \texttt{subsidy},
\texttt{penalty\_rate}) and objectives \texttt{-welfare.total},
\texttt{welfare.gini}. Firms maximize per-step profit against a stochastic
demand schedule; the regulator, government, and welfare layers implement
enforcement, redistribution, and aggregation. The Pareto trade-off is
aggregate welfare versus distributional equality across 32 governance scenarios.
Section~5.6 reports cross-domain validation at large scale.

\subsection{Wolf-Sheep Predation: Mean-Field Reimplementation}
\label{sec:wolfsheep}

To demonstrate compatibility with \emph{existing published models}, we
re-implement the canonical Mesa Wolf-Sheep Predation model~\cite{wilensky1997wolfsheep,kazil2020mesa}---the flagship example distributed with the Mesa library---as a mean-field ODE Arena.
Its published dynamics are Lotka-Volterra~\cite{lotka1925elements,volterra1926variazioni} predator-prey with grass regrowth on a
20$\times$20 spatial grid.

\textbf{HEAS reimplementation.} We implement the mean-field ODE approximation as a
4-layer Arena, deriving each term directly from Mesa's energy mechanics.
All six published parameters are used verbatim.

\textit{Grass.} Mesa's binary GrassPatch agents collapse to a continuous
density $G\in[0,1]$:
\begin{equation}
  \Delta G = \frac{1-G}{\tau_g} - \frac{S\,G\,\gamma}{N},
  \label{eq:grass}
\end{equation}
where $\tau_g\!=\!30$, $N\!=\!400$, and $\gamma\!\in\![0,1]$ is the policy
grazing-rate gene.

\textit{Sheep} ($g_s\!=\!4$, $r_s\!=\!0.04$):
\begin{equation}
  \Delta S = \Bigl(r_s - \frac{\max(0,\,1-g_s G\gamma)}{g_s}\Bigr)S
             \;-\; \frac{W\,S}{N}.
  \label{eq:sheep}
\end{equation}

\textit{Wolves} ($g_w\!=\!20$, $r_w\!=\!0.05$; $h$ = harvest-rate gene):
\begin{equation}
  \Delta W = r_w \frac{W\,S}{N}
             \;-\; \frac{\max(0,\,1-g_w S/N)}{g_w}\,W
             \;-\; h\,W.
  \label{eq:wolf}
\end{equation}

Every coefficient maps directly to Mesa defaults with no ad-hoc parameters. The
spatial grid is collapsed to a density field. Adding EA optimization to this
published model required zero additional coupling code: the objective function
reads \texttt{metrics\_episode()} keys, identical to the other two case studies.

A 30-run algorithm comparison on this model is reported in Section~5.5,
where the landscape structure determines which optimizer wins.

% ============================================================
\section{EVALUATION}

\subsection{Research Questions and Evaluation Design}

We define \emph{aggregation inconsistency} at two levels. \emph{Structural inconsistency} occurs when two pipeline stages implement logically distinct aggregation functions on the same \texttt{metrics\_episode()} key---for example, the optimizer reading final-step biomass while the tournament reads episode-mean biomass. \emph{Outcome inconsistency} is its observable consequence: the two stages produce non-identical policy rankings, measured as Kendall's~$\tau < 1$ between the orderings. The controlled experiment (\S\ref{sec:controlled-agg}) tests for outcome inconsistency as the operationalisation of structural divergence.

Our evaluation addresses two questions. First: does the \texttt{metrics\_episode()}
contract prevent aggregation inconsistency by construction? We test this through
a controlled experiment (\S5.2) in which we compare HEAS
against ad-hoc baselines that introduce deliberate aggregation heterogeneity across
pipeline stages. Second: is the contract portable across ODE-style and tabular
ABM models without rewriting framework logic? We test this through three case
studies spanning ecological, enterprise, and published-model reimplementation
domains (\S5.5).

The tournament stability results (\S5.3) and out-of-distribution
champion robustness (\S5.4) are \emph{secondary} evidence: they
validate that the contract enforces consistent semantics in practice, not that
it prevents inconsistency per se. \emph{Falsification} would require: (a) rank
reversals under HEAS despite contract enforcement, or (b) framework logic
modification required to port to a new case study.

A positive result on both questions would validate the core design hypothesis from~\S\ref{sec:intro-contributions}: that enforcing a single \texttt{metrics\_episode()} contract at the framework level prevents the ranking inconsistencies observed in ad-hoc coupling code.

\emph{Statistical reporting standards:} All confidence intervals use the BCa
bootstrap method (10,000 resamples) unless otherwise noted. Binomial proportions
use Wilson score intervals. Pairwise comparisons use Wilcoxon rank-sum tests.
Multiple comparisons apply Bonferroni correction.


\subsection{Controlled Aggregation Experiment: Metric Contract Effect Isolation}
\label{sec:controlled-agg}

For a metric contract to be useful, it must prevent the aggregation inconsistencies
that emerge when evolution, tournament, and inference stages read metrics differently.
We test this claim directly by implementing two pathways:

\begin{enumerate}
\item \emph{Uniform contract (HEAS)}: Single \texttt{metrics\_episode()} dict consumed
identically by evolutionary optimizer, tournament scorer, and statistical pipeline.
\item \emph{Ad-hoc aggregation (no contract)}: EA extracts \emph{final-step} biomass;
tournament reads \emph{episode-mean} biomass; statistics read \emph{rolling 10-step mean}
biomass. All other code identical; only metric aggregation differs.
\end{enumerate}

We run 15 independent NSGA-II search processes at small scale (150 steps, pop=20,
ngen=10) on the ecological case study, extract the Pareto front, and score each
candidate across 8 scenarios using four voting rules (argmax, majority, Borda,
Copeland). Rank reversals---cases where the tournament winner differs under uniform
vs.\ ad-hoc aggregation---measure contract efficacy.

Results from 15 runs with 8 scenarios and 4 voting rules per run:

\begin{table}[ht]
\centering\small
\caption{Aggregation Condition Comparison. A controlled aggregation experiment isolates
the effect of the metric contract by comparing HEAS (uniform interface) against ad-hoc
aggregation (independent per-stage implementations). Rank reversal = disagreement on
top-2 candidate between any two voting rules within a single run. Results confirm that the metric contract is the necessary mechanism for consistent aggregation: HEAS achieves 0\% rank reversals while ad-hoc pathways produce systematic divergence.}
\label{tab:agg-control}
\begin{tabular}{lccc}
\toprule
Condition & Rank Reversals & Voting Agreement & Pareto $\tau$ \\
& (out of 15 runs) & $P(\text{all agree})$ & vs.\ HEAS \\
\midrule
HEAS (uniform) & 0 (0.0\%) & 1.0 & 1.0 \\
Ad-hoc-Step & 1 (6.7\%) & 0.80 & 0.85 \\
Ad-hoc-Mean & 1 (6.7\%) & 0.93 & 0.92 \\
\bottomrule
\end{tabular}
\end{table}

Effect size (HEAS vs.\ ad-hoc): Cohen's $h = 0.476$ [95\% CI: 0.211--0.680],
$p < 0.05$ (binomial test, two-sided). The uniform metric contract eliminates
rank reversals entirely (0\%), while ad-hoc-Step aggregation produces reversals in 6.7\%
of runs ($\Delta = 6.7\%$, 95\% CI [1.2\%, 29.8\%]). Pareto front membership divergence,
measured by Kendall's $\tau$ between HEAS and ad-hoc fronts, shows that the two
pathways evolve qualitatively different solutions ($\tau = 1.0$ for HEAS vs.\ 0.85 for
Ad-hoc-Step). This result establishes the metric contract as a necessary mechanism for consistent
aggregation, not merely a coding convenience. In Stage 2 (stochastic robustness test), we
re-evaluated the same policies under a different random seed: HEAS maintained perfect
ranking consistency ($\tau = 1.000$), while both ad-hoc conditions also showed $\tau = 1.000$,
confirming that the deterministic ranking consistency emerges from the metric implementation
structure rather than external validation.

\subsection{Mesa Comparison: Coupling Code Reduction}
\label{sec:mesa-comparison}

We implement the ecological case study as a competent Mesa~3.3.1 model and
count non-blank, non-comment coupling code (model dynamics excluded).
Measured totals: 160~LOC (Mesa) vs.\ 5~LOC (HEAS)---97\% reduction.
To establish a fairer baseline, we additionally implement a reusable
utility-library abstraction (\texttt{mesa\_eco\_util.py}) that wraps
DataCollector setup, metric extraction, episode runners, and seed management
into a single \texttt{MesaEpisodeRunner} class. The call-site coupling code
for this utility-library approach is 42~LOC (measured)---still an 88\%
reduction compared to HEAS's 5~LOC. The utility class itself requires
$\sim$55~additional lines of one-time per-project infrastructure.

\begin{table}[ht]
\caption{Coupling code: Mesa~3.3.1 vs.\ HEAS vs.\ Mesa+utility-library.
  Mesa column: measured from \texttt{mesa\_eco.py}.
  Mesa+Util column: measured from \texttt{mesa\_eco\_util.py} (call-site LOC
  when using a reusable utility class; utility class itself is $\sim$55~LOC,
  a one-time per-project cost). HEAS column: measured from \texttt{heas/eco.py}.}
\label{tab:loc}
\centering\small
\begin{tabular}{p{0.44\linewidth}rrr}
\toprule
Task & Mesa & Mesa+Util & HEAS \\
\midrule
Metric contract / DataCollector setup   & 15 & 8  & 0 \\
Episode metric extraction               & 20 & 6  & 0 \\
Sequential episode runner               & 22 & 4  & 1 \\
Parallel runner (\texttt{ProcessPoolExecutor}) & 20 & 5  & 0 \\
EA fitness glue                         & 14 & 4  & 3 \\
Tournament scorer                       & 18 & 4  & 0 \\
Per-run seed management (30-run)        & 12 & 3  & 0 \\
Bootstrap CI                            & 35 & 6  & 0 \\
Adding a second objective (incremental) & 4  & 2  & 1 \\
\midrule
\textbf{Total}                          & \textbf{160} & \textbf{42} & \textbf{5} \\
\bottomrule
\end{tabular}
\end{table}

Table~\ref{tab:loc} shows the full comparison including the Mesa+utility-library
baseline (42~LOC, measured). Adding a metric requires editing DataCollector, EA
function, and tournament scorer in sync; in HEAS this is a one-line change in one
file. HEAS's claim over this baseline is that
\texttt{metrics\_episode()} is enforced at the \emph{framework} level, so new
projects inherit the uniform contract without writing or maintaining a
utility module.

\subsection{Multi-Scale Reproducibility}
\label{sec:multiscale}

For RQ2, we test whether the ecological study is reproducible across the
small- and large-scale settings used in this paper. At small scale
($n\!=\!30$, pop\,=\,20, ngen\,=\,10, steps\,=\,150),
ecological HV\,=\,7.665\,(std\,3.518, CI\,[6.424,\,8.914]; HV reference point:
mean\_biomass\,=\,$-$150, CV\,=\,0.5).
The distribution is bimodal, and a $4\!\times\!4$ budget ablation
($n\!=\!20$ runs/cell; S2c) identifies population size as the controlling
factor: escape rate (HV\,$>\!9.0$) rises from 15--25\% at pop\,$\le\!20$ to
40--80\% at pop\,$\ge\!30$.

The study was then replicated at larger scale
(steps\,=\,1,000, n\_eval\,=\,5 episodes per genome evaluation, pop\,=\,30,
ngen\,=\,15, $n\!=\!20$ runs). Ablation (S2c) identifies population size as
controlling factor. At this horizon NSGA-II outperforms random search: mean HV\,=\,16.66\,(std\,10.84,\,CI\,[12.37,\,20.94])
versus random mean\,=\,6.00\,(std\,0.34,\,CI\,[5.85,\,6.14]).
Wilcoxon test: $p\!=\!1.91\!\times\!10^{-6}$; Cohen's $d\!=\!1.39$
(large effect); NSGA-II achieves 178\% higher mean HV than random search.
At 150-step scale the same comparison is not significant ($d\!=\!0.49$,
$p\!=\!0.49$); the Pareto landscape becomes meaningfully richer only at the
longer horizon.

\begin{figure}[ht]
  \centering
  \includegraphics[width=\columnwidth]{fig4_large_scale_showdown}
  \caption{Large-scale algorithm showdown: NSGA-II vs.\ Random at 1,000
    steps/episode ($n\!=\!20$ runs). Bars show mean\,$\pm$\,1\,std; error bars
    show 95\% CI; dots show individual runs.
    Cohen's $d\!=\!1.39$; Wilcoxon $p\!=\!1.91\!\times\!10^{-6}$.}
  \label{fig:large-scale}
\end{figure}

Figure~\ref{fig:large-scale} shows per-run distributions and aggregate
statistics. The HEAS framework required no code changes between the 150-step
and 1,000-step studies: only the \texttt{steps} argument changed.

\subsection{Tournament Validation and OOD Champion Robustness}
\label{sec:tournament}

For RQ3, we evaluate stable tournament behavior and champion performance in the
tested ecological scenario families. Voting-rule agreement and sample complexity
were assessed across 30 repeats of
the 8-scenario tournament with 100 episodes per scenario. All four voting rules
agreed in 100\% of
(scenario, repeat) pairs, and $P(\text{correct winner})\!=\!1.0$ at all tested
episode budgets (4, 10, 25, 50, 100 per scenario). The 100\% agreement likely
reflects clear dominance on these low-dimensional problems: with only 2--4 policy
genes and 8 scenarios, the Pareto landscape offers unambiguous winners. A
near-tie sensitivity study~(S9) confirms that this unanimity disappears when
candidate policies have statistically indistinguishable means.

Tournament rankings remain stable to measurement noise up to 0.65\% of the inter-policy margin (details in S3).

OOD champion robustness was validated at two scales. At small scale (16
scenarios, steps\,=\,500), the champion won all 16 ($p\!<\!0.001$). At large
scale (32 scenarios), the champion policy
won all 32 out of 32 ($p\!=\!4.66\!\times\!10^{-10}$, one-sided binomial, $H_0\colon \pi\!=\!0.5$ vs.\ $H_a\colon \pi\!>\!0.5$;
Wilson 95\% CI for win fraction: [0.94, 1.0]; $\Delta\!=\!+25.0$, $+3.96\%$;
$d\!=\!0.17$). The small effect size ($d\!=\!0.17$) reflects consistently tight score distributions across OOD scenarios: the champion generalises without overfitting to any single case, consistent with what we expect when policy evaluation is decoupled from training conditions. Win~rate (32/32) is the primary inferential statistic; paired deltas confirm that all individual scenario margins are positive. For comparison, a random policy drawn from the evolutionary search space wins approximately 50\% of scenarios under the same tournament; the champion's 100\% win fraction thus represents a substantial departure from random performance. In the single-anchor \emph{true OOD} study (Exp.~C, S6c), the
champion won all 24 scenarios where all four factors lay outside training
($p\!=\!1.19\!\times\!10^{-7}$), matching the 100\% in-distribution win rate.
Here win rate and paired deltas are more informative than Cohen's~$d$ because score distributions are near-deterministic; crucially, ranking does not reverse when environmental conditions shift outside the training distribution---validating the decoupling premise from~\S\ref{sec:intro-contributions}: a single \texttt{metrics\_episode()} implementation shared across all pipeline stages prevents the aggregation drift that would otherwise confound OOD evaluation.


\subsection{Algorithm Agnosticism and Scale Sensitivity}
\label{sec:scale}

For RQ4, we test whether framework usefulness depends on any single optimizer.
It does not. In the ecological ablation at steps\,=\,300
(pop\,=\,20, ngen\,=\,10, $n\!=\!10$ runs), Simple hillclimbing outperforms
NSGA-II (Simple\,=\,19.66, NSGA-II\,=\,9.99, Random\,=\,7.61). On this
2-gene landscape, the trade-off is narrow and local search is sufficient. HEAS
therefore remains useful even when the best optimizer is not NSGA-II.

Scale sensitivity is shown in Figure~\ref{fig:heatmap}. Mean HV generally rises
with longer simulation horizons, from 6.0 at 140 steps and 5 episodes per
evaluation to 10.1 at 300 steps and 20 episodes. This indicates that richer
episode horizons produce a more informative search landscape.

\begin{figure}[ht]
  \centering
  \includegraphics[width=0.846\columnwidth]{fig3_scale_heatmap}
  \caption{Scale sensitivity heatmap: mean HV by simulation horizon
    (steps\,$\times$\,episodes/genome). Each cell is $n\!=\!10$ independent
    NSGA-II runs. HV generally increases with steps; non-monotonicities
    between the 300- and 500-step rows reflect high between-run variance.
    Episode count has a weaker secondary effect.}
  \label{fig:heatmap}
\end{figure}

Given the 9 cells in the $3\times 3$ grid, we applied Bonferroni correction
($\alpha = 0.05 / 9 = 0.0056$) to pairwise comparisons across budget conditions.
No cells remained significantly different after correction, confirming that HV
variance across budget conditions is primarily stochastic rather than systematic.


The Wolf-Sheep comparison reaches the same practical conclusion from a different
direction. At steps\,=\,200 (pop\,=\,50, ngen\,=\,25, $n\!=\!30$ runs), Random
achieves the highest HV (14.35), followed by NSGA-II (10.36) and Simple (9.14);
Random also significantly exceeds NSGA-II
($p\!=\!7.99\!\times\!10^{-6}$). The point is not that one optimizer wins in
general. The point is that optimizer performance is landscape-dependent, and
HEAS makes that comparison possible without rewriting framework code.

\subsection{Cross-Domain Portability}
\label{sec:crossdomain}

For RQ5, we ask whether the same framework code can be reused
across domains. The answer is yes at the code level. The ecological and
enterprise studies use the same optimization, tournament, and statistical
infrastructure; only the domain-specific \texttt{metrics\_episode()}
implementations differ. In the enterprise tournament validation across 8
governance scenarios (20 repeats, 30 episodes per scenario), all four voting
rules agree in 100\% of cases and $P(\text{correct winner})\!=\!1.0$ at all
tested budgets. This supports a narrow portability claim: the framework
transfers across the tested ODE/tabular domains without rewriting framework
logic. It does not yet establish portability to spatially-explicit or
heterogeneous-agent ABMs.

\subsection{Threats to Validity}

\textbf{Construct validity.} The case studies are composition vehicles, not
calibrated domain models; results are evidence about framework behavior. The
metric contract prevents \emph{structural divergence across consumers}: all
three pipeline stages read the same implementation. An incorrect \texttt{metrics\_episode()}
implementation will be consistently wrong rather than silently inconsistent. This framework prevents structural inconsistency (different aggregation implementations), but does not enforce semantic consistency (correct interpretation of metric values by developers). The
framework additionally exposes \texttt{validate\_metrics\_episode()}, which raises
\texttt{TypeError} when a \texttt{metrics\_episode()} implementation returns a
non-dict or non-float value, catching contract violations early in development.

\textbf{Internal validity.} Fixed evaluation seeds isolate framework effects
but can suppress variance in some metrics; zero-variance outcomes are treated as
diagnostic signals (e.g., the $K\!=\!120$ degeneracy, enterprise determinism)
rather than performance claims. OOD validation at steps\,=\,500 and on
single-anchor extrapolation scenarios (Exp.~C) extends evaluation beyond the
training horizon, but only within the tested family of ecological scenarios.

\textbf{External validity.} HEAS targets ODE/tabular dynamics; spatially-explicit ABMs
with heterogeneous agents on grids lie outside the current scope.

% ============================================================
\section{CONCLUSION}

The paper supports a focused claim: HEAS makes simulation-based policy search
more auditable and more reusable by enforcing one metric contract across
optimization, comparison, and inference. The evidence is consistent on this
point. HEAS reduces coupling code sharply relative to Mesa-based baselines,
supports reuse across three case studies in the targeted model family, and
keeps tournament and inference stages on the same metric pathway used in search.

The results also clarify scope. HEAS is not a claim about universal optimizer
superiority. In richer ecological settings, NSGA-II performs strongly. In
simpler landscapes, other search strategies perform as well or better. That is
a feature of the framework rather than a weakness: optimizer choice can vary
without changing framework logic.

\textbf{Limitations.} HEAS targets ODE/tabular dynamics, not spatial ABMs. Its scalar
metric contract (\texttt{dict[str, float]}) excludes time-series outputs.

An anonymized artifact package containing all scripts and result files (S1--S10)
is available for reviewers; upon acceptance, source code and raw data will be
published with DOI.

% ============================================================
\appendix
\section{Reported Means and Sources}

\begin{table}[ht]
\caption{Principal results reported in this paper with source file IDs (S1--S10).}
\label{tab:appendix-results}
\centering\small
\begin{tabular}{p{0.52\linewidth}p{0.23\linewidth}p{0.1\linewidth}}
\toprule
Quantity & Value & Source \\
\midrule
Coupling LOC reduction & 97\% (160 vs.\ 5) & S1 \\
NSGA-II vs.\ Random: large scale HV & $d=1.39$, $p<10^{-6}$ & S5 \\
32-scenario OOD champion wins & 32/32, $p=4.66\!\times\!10^{-10}$ & S6b \\
True OOD 24-scenario wins & 24/24, $p=1.19\!\times\!10^{-7}$ & S6c \\
Enterprise portability & Validated (tournament 100\%) & S8 \\
\bottomrule
\end{tabular}
\end{table}

\noindent\textbf{Source Files.} All experiment data and scripts (S1--S8, supplementary ablations)
are provided in the artifact package under \texttt{experiments/results/}.

% ============================================================
\begingroup
\small
\bibliographystyle{wsc}
\bibliography{references}
\endgroup

% ============================================================
\section*{AUTHOR BIOGRAPHIES}

\noindent {\bf \MakeUppercase{Ruiyu Zhang}} is affiliated with the Department of Politics and Public Administration at The University of Hong Kong, Hong Kong SAR, China. Her e-mail address is \email{Ruiyuzh@connect.hku.hk}.\\

\noindent {\bf \MakeUppercase{Lin Nie}} is affiliated with the Department of Applied Social Sciences at The Hong Kong Polytechnic University, Hong Kong SAR, China. His e-mail address is \email{lin-apss.nie@polyu.edu.hk}.\\

\noindent {\bf \MakeUppercase{Xin Zhao}} is affiliated with the Department of Applied Social Sciences at The Hong Kong Polytechnic University, Hong Kong SAR, China. His e-mail address is \email{xinnn.zhao@connect.polyu.hk}.

\end{document}
